\section{Análisis de resultados}
\label{sec:analisis}

Tras ejecutar los experimentos y almacenar los resultados en la base de datos SQLite, el siguiente paso es procesar esta información para obtener gráficos y estadísticas comparativas. El proyecto proporciona un sistema modular y automatizado de análisis post-ejecución.

\subsection{Estructura del script de análisis}

El análisis de resultados está completamente centralizado a través del script:
\begin{verbatim}
scripts/analisis.py
\end{verbatim}
Este script actúa como punto de entrada de la biblioteca modular del análisis (\texttt{scripts/analisis\_pkg/}), la cual está compuesta por:
\begin{itemize}
    \item \texttt{cache.py}: Gestión de la memoria caché al consultar datos de SQLite para acelerar el procesamiento.
    \item \texttt{config.py}: Carga las configuraciones de directorios, experimentos y el propio análisis.
    \item \texttt{parser.py}: Filtra y valida los nombres de archivos CSV y datos de rendimiento.
    \item \texttt{plotting.py}: Lógica gráfica para boxplots, violinplots y curvas de convergencia.
    \item \texttt{processing.py}: Motor principal que organiza el procesamiento en lotes paralelos usando Joblib para evitar el desbordamiento de memoria RAM.
\end{itemize}

El análisis consulta los experimentos finalizados directamente desde la base de datos de SQLite, procesando las iteraciones (contenidas en formato CSV como BLOBs) de forma asíncrona y robusta.

\subsection{Configuración del análisis}
\label{subsec:config_analisis_v2}

El proceso de análisis se configura mediante archivos JSON ubicados en \texttt{src/util/json/}:

\begin{itemize}
    \item \texttt{dir.json}: Especifica las rutas relativas donde se guardan los resultados finales, incluyendo resúmenes de rendimiento (\texttt{"resumen"}), gráficos de convergencia (\texttt{"graficos"}), boxplots (\texttt{"boxplot"}), violinplots (\texttt{"violinplot"}), entre otros.
    
    \item \texttt{analysis.json}: Este archivo actúa como panel de control para el rendimiento del procesamiento y la visualización gráfica. Su estructura es la siguiente:
    \begin{quote}
    \begin{verbatim}
{
    "performance": {
        "batch_size": 200,      // Tamaño del lote a procesar
        "gc_interval": 500,     // Frecuencia del Garbage Collector
        "ram_warning": 85,      // % de uso de RAM para advertencia
        "cache_max_size": 256   // Tamaño máximo de caché en base de datos
    },
    "graficos": {
        "graficos_por_corrida": false, // Evita generar miles de gráficos individuales
        "modo_logaritmico": "auto",    // Escala ('log_scale', 'log_transform', 'none', 'auto')
        "dpi": 300                     // Resolución de gráficos generados
    }
}
    \end{verbatim}
    \end{quote}
    
    \item \texttt{experiments\_config.json}: El script de análisis consulta este archivo para obtener la lista de metaheurísticas activas (\texttt{"mhs"}) e instancias seleccionadas a incluir en las tablas estadísticas comparativas.
\end{itemize}

\subsection{Ejecución del análisis}

Para procesar y analizar los resultados obtenidos de la base de datos local:

1. Asegúrese de haber completado y finalizado ejecuciones de experimentos.
2. Desde la terminal en el directorio raíz del proyecto, ejecute:
    \begin{verbatim}
    python scripts/analisis.py
    \end{verbatim}

El script procesará los datos pendientes de la base de datos en lotes, optimizando la memoria RAM, e informando el progreso en tiempo real.

\subsection{Salidas generadas}

Los resultados procesados se escriben en el directorio \texttt{outputs/results/}, organizados en:
\begin{itemize}
    \item \textbf{Tablas estadísticas (\texttt{resumen/})}:
    \begin{itemize}
        \item Archivos CSV con métricas detalladas por metaheurística (mejor fitness, promedio de fitness, desviación estándar, tiempo promedio de ejecución).
        \item Estadísticas de exploración y explotación promedio (\%XPL y \%XPT) a lo largo de las iteraciones.
    \end{itemize}
    
    \item \textbf{Visualizaciones gráficas}:
    \begin{itemize}
        \item \textbf{\texttt{graficos/}}: Historial de búsqueda de mejores soluciones (fitness vs iteraciones), y gráficos representativos de balance \%XPL y \%XPT por algoritmo.
        \item \textbf{\texttt{boxplot/}} y \textbf{\texttt{violinplot/}}: Distribución de calidad final de las corridas para cada algoritmo, permitiendo comparar visualmente la consistencia y robustez de los optimizadores.
    \end{itemize}
\end{itemize}
% --- Fin de la Sección ---